% This document was generated by an LLM.

\documentclass{essay}

\title{The Digits Are Not the Number}
\author{}
\date{}

\begin{document}

\maketitle

\section{A picture and its immediate failure}

Ask someone what an irrational number is and you will often be told: a number
whose decimal expansion goes on forever. The picture behind the answer is a
tape of digits unrolling to the right without end, never settling down.

The picture is not so much wrong as uninformative, and it is worth being
precise about why. Consider
\[
  \tfrac13 = 0.333\ldots, \qquad \tfrac12 = 0.5000\ldots, \qquad
  \sqrt2 = 1.41421356\ldots
\]
All three expansions are infinite. Under the base-$10$ algorithm every real
number receives an infinite string of digits; terminating expansions are those
that terminate in an infinite tail of zeros. So infinitude of digits is a
property of the algorithm, shared by every real, and it separates nothing.

What actually separates the cases is periodicity.

\begin{proposition}\label{prop:periodic}
A real number $x$ is rational if and only if its base-$b$ expansion is
eventually periodic, for any (equivalently, every) integer base $b \ge 2$.
\end{proposition}

The proof in one direction is long division: computing the expansion of $p/q$
produces at each step a remainder in $\{0,1,\ldots,q-1\}$, and the state of the
algorithm is determined by that remainder, so within $q$ steps some remainder
recurs and the process cycles. The converse is the geometric series: an
eventually periodic expansion is a finite decimal plus a scaled geometric sum
with rational ratio $b^{-k}$, hence rational.

Proposition~\ref{prop:periodic} is a genuine theorem, and it does exactly what
the naive picture failed to do. It also has an instructive shape. Rationality
is a base-independent notion --- $p/q$ is a ratio of integers regardless of how
we write it --- and the proposition asserts that eventual periodicity, which
looks base-dependent, is in fact base-independent too. That the equivalence
holds in every base is a fact requiring proof, not a triviality.

But notice what has happened. We began with a proposed definition of
irrationality and ended with a characterization by a property of a
representation, a property whose base-independence is itself a theorem. That
inversion is the subject of this essay. The digit expansion is a coordinate
system imposed on the real line, and like any coordinate system it has features
of its own that do not belong to the object it describes. Below I argue that
the decimal representation is in three specific respects a poor coordinate
system; that better ones exist; that the difference between them is not
cosmetic but reveals structure invisible from the digit side; and that pushing
on the question of what a representation must supply leads, fairly quickly, to
questions about what it means for a real number to exist at all.

\section{Three defects}

\subsection{Base dependence}

The number $1/3$ has the base-$10$ expansion $0.333\ldots$ and the base-$3$
expansion $0.1$. The number $1/2$ terminates in base $10$ and repeats in base
$3$, where it is $0.111\ldots$. Nothing about the number changes; the visible
texture of its representation changes completely.

Termination, then, is not a property of a real number. It is a property of a
pair: a number and a base. A rational $p/q$ in lowest terms terminates in base
$b$ exactly when every prime dividing $q$ divides $b$. The distinction the
decimal system makes most salient --- between $0.5$, which looks finite and
tame, and $0.333\ldots$, which looks infinite and unruly --- is an artifact of
the factorization of $10$.

Proposition~\ref{prop:periodic} shows that one can recover a base-independent
notion (rationality) from a base-dependent formalism (periodicity), but only by
proving that the base does not matter. Any statement about digits carries this
burden. Until the burden is discharged, a digit-level observation is a
statement about notation.

\subsection{Non-uniqueness}

The map from digit strings to real numbers is not injective:
\[
  0.999\ldots = 1.000\ldots
\]
The identity is forced by the definition of the expansion as the sum of a
series: $\sum_{n\ge1} 9 \cdot 10^{-n} = 1$. Every terminating expansion has a
second representation ending in repeated $(b-1)$'s, and these are the only
collisions.

The consequence for our purposes is structural rather than arithmetical. If a
real number simply \emph{were} its digit string, then $0.999\ldots$ and $1$
would be different numbers, which they are not. To make the identification work
one must first quotient the space of digit strings by the relation just
described. The corrected slogan --- a real number is an equivalence class of
digit strings --- is already conceding the point: the digit string is not the
number, because knowing the number does not determine the string.

I will not belabor this. It is the best-known item in this catalogue and the
shallowest.

\subsection{Carrying, or: the representation is arithmetically inert}

The serious defect is that one cannot compute with digit expansions.

Consider adding $x = 0.333\ldots$ and $y = 0.666\ldots$. Suppose an adversary
hands you the digits one at a time and asks for the first digit of $x+y$ after
you have seen any finite number of them. After a million $3$'s and a million
$6$'s the sum of the visible parts is $0.999\ldots9$, and you must announce
either $0$ or $1$ as the leading digit of the result. You cannot. If the
expansions continue as $3$'s and $6$'s forever, the sum is exactly $1$ and the
answer is $1$; if the very next digit of $x$ is $2$ instead of $3$, the sum
begins $0.9$ and the answer is $0$. No finite prefix of the inputs determines
even the first digit of the output.

The phenomenon is unbounded carry propagation, and it is not an edge case
cleverly constructed for the occasion. It is generic near any point with two
expansions, that is, near every terminating rational, a dense set. State the
consequence precisely:

\begin{proposition}
Addition is not computable as a function from pairs of digit streams to digit
streams. There is no algorithm which, reading the digits of $x$ and $y$ in
order, emits the digits of $x+y$ in order, with each emitted digit depending on
only finitely many read digits.
\end{proposition}

This is a strong statement and it disqualifies the representation for any
purpose involving actual computation. The field operations, the operations that
make $\R$ what it is, do not act on digit streams in any local way. A
representation on which one cannot add is not a representation of a field; it
is a labelling scheme.

That the defect was recognized early is worth recording. Turing's 1936 paper
was called ``On Computable Numbers,'' and he defined a computable real as one
whose decimal expansion is produced by a machine. The correction he published
the following year was concerned precisely with the inadequacy of that choice.
Modern computable analysis does not use digit streams. It represents a real by
a sequence of nested rational intervals of shrinking width, or by a Cauchy
sequence of rationals together with an explicit modulus of convergence:
\[
  (q_n)_{n \in \N} \subset \Q
  \quad \text{with} \quad
  |q_n - q_m| \le 2^{-n} \ \text{ for } m \ge n .
\]
On these representations the field operations are computable, because an
interval enclosing $x$ and an interval enclosing $y$ combine to give an
interval enclosing $x+y$, and the enclosures can be made as tight as desired by
going deeper into the inputs. The information is the same; its organization is
not.

Here is a test for the naive picture, worth carrying around. What is the digit
expansion of
\[
  \sqrt2 + \bigl(2 - \sqrt2\bigr) ?
\]
The answer is $2$, and the reason is the field axioms. No digit-level procedure
reaches it: both summands have expansions that are, so far as anyone can prove,
statistically structureless, and the cancellation that makes the answer exact
is invisible at every finite stage. The gap between what you know and what the
digits show you is the measure of how much the representation is not the
number.

\section{What the object is}

If the digits are a chart, what is being charted?

There are several standard constructions and they agree. Dedekind's takes a
real number to be a cut: a partition of $\Q$ into a downward-closed set $A$
with no greatest element and its complement $B$. The irrationals are the cuts
whose associated $A$ has no rational supremum; $\sqrt2$ is the cut
$A = \{q \in \Q : q < 0 \text{ or } q^2 < 2\}$, and the construction says
nothing about digits. Cantor's takes a real to be an equivalence class of
Cauchy sequences of rationals under the relation $(a_n) \sim (b_n)$ iff
$a_n - b_n \to 0$; the irrationals are the classes containing no constant
sequence.

The axiomatic view subsumes both.

\begin{theorem}
There is, up to isomorphism of ordered fields, exactly one complete ordered
field.
\end{theorem}

$\R$ is that field. Both constructions produce it; both are therefore
descriptions of the same object, and neither is more the real numbers than the
other. What matters for the present argument is the negative observation: no
formulation of the definition mentions digits, or bases, or expansions. The
decimal system is an afterthought, a convenience introduced once the object is
in hand, and every property one might read off a digit string must be
reconnected to the definition before it counts as a property of a number.

\section{A better representation}

Criticism of a representation is cheap unless one can exhibit an alternative.
The continued fraction expansion is the alternative, and comparing the two
converts a complaint into an argument.

Every real $x$ has an expansion
\[
  x = a_0 + \cfrac{1}{a_1 + \cfrac{1}{a_2 + \cfrac{1}{a_3 + \dotsb}}},
  \qquad a_0 \in \Z, \quad a_i \in \N_{\ge 1} \ (i \ge 1),
\]
written $x = [a_0; a_1, a_2, a_3, \ldots]$. It is produced by the obvious
algorithm: take the integer part, invert the remainder, repeat. For irrational
$x$ the expansion is infinite and unique.

Four points of comparison.

\medskip
\noindent\textbf{It is base-free.} The algorithm involves no choice of $b$. The
first defect vanishes because there is nothing for it to attach to.

\medskip
\noindent\textbf{The rational/irrational split is termination, not
periodicity.} The algorithm applied to $p/q$ is the Euclidean algorithm and
halts; applied to an irrational it does not. So the dichotomy that decimals
express through the subtle condition ``eventually periodic'' --- a condition
requiring Proposition~\ref{prop:periodic} to connect it to anything --- appears
here as the crudest possible one.

\medskip
\noindent\textbf{Periodicity is then free to mean something.} It does.

\begin{theorem}[Lagrange]
An irrational $x$ has an eventually periodic continued fraction expansion if
and only if $x$ is a quadratic irrational, that is, a root of a quadratic with
integer coefficients.
\end{theorem}

Thus $\sqrt2 = [1;2,2,2,\ldots]$ and
$\tfrac{1+\sqrt5}{2} = [1;1,1,1,\ldots]$. In the decimal system, periodicity is
spent on distinguishing $1/7$ from $1/2$, a distinction of no structural
interest; here it identifies the first nontrivial layer of the algebraic
hierarchy. The same formal feature, allocated to a different representation,
detects a different and much better thing.

\medskip
\noindent\textbf{The truncations are optimal.} The convergents $p_n/q_n$
obtained by cutting the expansion at $a_n$ satisfy
\[
  \left| x - \frac{p_n}{q_n} \right| < \frac{1}{q_n q_{n+1}} \le
  \frac{1}{q_n^2},
\]
and, more to the point, they are best approximations in a strong sense: if
$q \le q_n$ and $p/q \ne p_n/q_n$ then $|qx - p| > |q_n x - p_n|$. Every
exceptionally good rational approximation to $x$ is a convergent. Decimal
truncation offers nothing comparable; $3.14159$ approximates $\pi$ to six
figures with denominator $10^5$, while the convergent $355/113$ does better
with denominator $113$.

The comparison establishes the essay's thesis in its strong form. The two
representations carry the same information --- each determines the number,
hence the other --- and yet one makes the arithmetic of approximation legible
and the other conceals it. Legibility is not a property of the number. It is a
property of the chart.

\section{What the good chart reveals}

Continued fractions do more than tidy up. They open onto a quantitative theory
of irrationality that the digit picture cannot see at all.

The question is: how well can a given irrational be approximated by rationals
with small denominators? Define the irrationality measure $\mu(x)$ as the
infimum of exponents $\mu$ such that
\[
  \left| x - \frac{p}{q} \right| < \frac{1}{q^{\mu}}
\]
has only finitely many rational solutions $p/q$. This is base-independent by
construction, and it stratifies the irrationals.

Dirichlet's theorem gives $\mu(x) \ge 2$ for every irrational: the inequality
with exponent $2$ always has infinitely many solutions, the convergents among
them. At the other end, Liouville showed that an algebraic number of degree $d$
cannot be approximated to order better than $d$, and used this to construct the
first numbers proved transcendental --- the point being that a number
approximated \emph{too} well by rationals cannot be algebraic. Liouville's
constant
\[
  L = \sum_{k=1}^{\infty} 10^{-k!} = 0.110001000000000000000001\ldots
\]
has $\mu(L) = \infty$. Roth's theorem (1955) closed the gap at the top: every
irrational algebraic number has $\mu = 2$ exactly, the minimum. Algebraic
irrationals are, uniformly and regardless of degree, as badly approximable as
an irrational is permitted to be.

So there is a real hierarchy --- rational, quadratic irrational, higher
algebraic, transcendental, Liouville --- graded by approximability, and it is
detected by the growth of the partial quotients $a_n$: bounded partial
quotients mean badly approximable, huge partial quotients mean the convergents
are exceptionally good. The golden ratio, all of whose partial quotients equal
$1$, is the worst-approximable number there is, and this is the precise sense in
which it is the ``most irrational.'' None of this structure is visible in a
decimal expansion.

Fairness requires one concession. $L$ was written above as a decimal, and its
structure is entirely obvious there --- a sparse pattern of $1$'s at positions
$k!$. Liouville numbers are visible in the digit picture precisely because
extremely good rational approximation corresponds to extremely long runs of
zeros. The decimal representation is not uniformly blind; it is blind in a
specific direction. It sees sparseness and it cannot see approximation quality
in general.

\section{Almost every number, and no number in particular}

Push the digit picture in the one direction where it seems to promise a theory
--- the statistics of the digits themselves --- and it produces a result that is
worth examining closely, because the shape of its failure is instructive.

\begin{definition}
$x$ is normal in base $b$ if for every $k$, each of the $b^k$ possible blocks
of $k$ digits occurs in the base-$b$ expansion of $x$ with asymptotic frequency
$b^{-k}$. It is absolutely normal if it is normal in every base $b \ge 2$.
\end{definition}

\begin{theorem}[Borel, 1909]
Almost every real number, with respect to Lebesgue measure, is absolutely
normal.
\end{theorem}

The exceptional set has measure zero. And yet: it is not known whether $\pi$ is
normal in any base. Nor $e$, nor $\sqrt2$, nor $\log 2$, nor any constant that
arose from mathematics rather than being constructed for this purpose. Not one.
The numbers we can prove normal are the ones built to be normal, such as
Champernowne's constant
\[
  C_{10} = 0.123456789101112131415\ldots,
\]
formed by concatenating the integers, which is normal in base $10$ by
construction and was published as the first explicit example.

Two lessons. The first is about the digit picture: it supports a statistical
theory that proves almost everything about almost all numbers and almost
nothing about any number one can name. Whatever is doing the work in Borel's
theorem is measure theory, not the expansion; the expansion contributes the
statement and no leverage on it.

The second lesson is the more disquieting. ``Almost every real is normal, and no
named real is known to be'' is an instance of a pattern. The named reals form a
countable set --- there are only countably many finite descriptions --- and a
countable set has measure zero. The typical real, in the measure-theoretic
sense, is one we can never name, discuss, or single out. Every real number that
will ever appear in a paper lies in a null set. The continuum is mostly
inaccessible, and this is not a limitation of technique but a cardinality fact.

Which raises the question of what such numbers are, and whether a representation
should be required to hand us anything at all.

\section{Existence, and the point of a representation}

Turing's diagonal argument settles the accounting. There are countably many
Turing machines and uncountably many reals; therefore all but countably many
reals are uncomputable, in the sense that no algorithm produces arbitrarily good
rational approximations to them. Almost every real, again in the measure sense,
is a number whose digits --- or convergents, or nested intervals; the
representation is irrelevant here --- cannot be generated by any procedure
whatsoever.

The completeness of $\R$, which is the axiom that distinguishes it from $\Q$
and the engine of every limit argument in analysis, is what manufactures these
objects, and it does so non-effectively.

\begin{theorem}[Specker]
There is a computable, strictly increasing, bounded sequence of rationals whose
limit is not a computable real number.
\end{theorem}

The sequence is explicit: enumerate the halting Turing machines and sum
$2^{-n}$ over those that halt. The partial sums are computable and increase to a
bound; the limit encodes the halting problem. Monotone convergence, which every
analysis course treats as elementary, here converts a completely effective input
into a non-effective output. The supremum exists by completeness and cannot be
computed by anything.

One may accept this and move on, and classical mathematics does. But one may
also ask what it means to assert the existence of an object that no procedure
can approach, and the constructive tradition takes the question seriously. If a
real number is required to come equipped with a procedure --- if to have a real
number is to have a way of producing rational approximations to any specified
accuracy --- then the classical picture changes in ways that are not
adjustments at the margin.

Trichotomy fails. Given a constructive real $x$ one cannot in general decide
which of $x < 0$, $x = 0$, $x > 0$ holds: deciding it would require knowing
whether all the approximations that have not yet been computed will remain
within a shrinking interval around zero, and that is not finitely checkable.
Deciding it in general would settle the halting problem. The trichotomy law,
which looks like a triviality about the ordering of the line, turns out to
encode an assumption about what we are allowed to know.

Bishop's development of constructive analysis shows how much survives, which is
a surprising amount: continuous functions on closed intervals, differentiation,
integration, the essential machinery. What is lost is the freedom to case-split
on undecidable comparisons, and a handful of theorems that were secretly using
that freedom.

The point is not to endorse this position. The point is that the naive picture
we started with --- a real number is a tape of digits --- is not a neutral
placeholder awaiting refinement. It is a specific and contestable ontological
commitment. It says that a number is a completed infinite object, given all at
once, whether or not anything could ever produce it. That is exactly the
commitment the constructivist declines, and exactly the one that generates
Specker's sequence and the uncountable ocean of numbers nobody will ever name.
The tape of digits is not a chart of the reals. It is a chart of a particular
answer to the question of what a real number is.

\section{Why the bad chart survives}

Given all this, one should ask why the digit representation persists, and the
answer is that it is genuinely good at one thing.

Cantor's diagonal argument requires a representation in which the numbers are
laid out as sequences of independently specifiable symbols from a finite
alphabet, so that a new sequence can be constructed by disagreeing with the
$n$-th listed sequence in the $n$-th place. Decimal expansions provide this. The
proof of the uncountability of $\R$ --- the fact underlying every observation
above about how much of the line is inaccessible --- is a proof about digit
strings, and it needs no more from them than that they exist and nearly
determine the number. The $0.999\ldots$ ambiguity intrudes and is disposed of by
avoiding the digits $0$ and $9$ in the constructed number.

So the representation that cannot support addition is the one that establishes
the cardinality of the continuum. This is not paradoxical. Digit expansions
encode the reals as an infinite product of finite sets, and infinite products of
finite sets are exactly what combinatorial and set-theoretic arguments handle
well. What such an encoding discards is the field structure and the metric
structure --- which is to say, everything analysis is about. Use it for
cardinality and for measure-theoretic statements about typical numbers, where it
is the right tool. Do not use it for arithmetic, for approximation, or as a
definition.

\section{Conclusion}

Is it fair to think of irrational numbers as infinite strings of digits after
the dot? As a mnemonic, and with periodicity supplied in place of infinitude, it
is serviceable. As an account of what an irrational number is, it fails, and the
manner of the failure is more interesting than the fact.

It fails first on its own terms, being base-dependent, non-injective, and
incapable of supporting addition. It fails by comparison, since the continued
fraction expansion carries the same information while making visible a whole
theory --- approximability, the algebraic hierarchy, Lagrange's characterization
of quadratic irrationals --- that the decimal expansion conceals. And it fails
philosophically, in that it presupposes an answer to a live question, treating a
real number as a completed infinite object handed over in full when the great
majority of such objects can never be handed over at all.

The corrective is not a better representation. It is the recognition that a
representation is a choice, that each choice illuminates the structure it was
built to reflect and hides the rest, and that the number is the thing that
persists across all of them. $\sqrt2$ is the positive root of $x^2 = 2$: the
cut in $\Q$ at that point, the length of the diagonal of the unit square, the
unique real whose square is $2$. It is $1.41421356\ldots$ and it is
$[1;2,2,2,\ldots]$, and neither string is the number. They are shadows cast by
particular projections, and the mathematics is in the object, not in the
shadows.

\end{document}
